% SPDX-License-Identifier: CC-BY-NC-ND-4.0
% Copyright 2026 Ingolf Lohmann.
\documentclass[11pt,a4paper]{article}

\usepackage{fontspec}
\setmainfont{Latin Modern Roman}
\setsansfont{Latin Modern Sans}
\setmonofont{Latin Modern Mono}[Scale=0.88]
\usepackage[provide=*,main=ngerman,english]{babel}
\usepackage{microtype}
\usepackage{geometry}
\usepackage{amsmath,amssymb,mathtools,amsthm}
\usepackage{booktabs,longtable,tabularx,array}
\usepackage{enumitem}
\usepackage{xcolor}
\usepackage{fancyhdr}
\usepackage{seqsplit}
\usepackage{xurl}
\usepackage{hyperref}
\usepackage{bookmark}
\usepackage[nameinlink,noabbrev]{cleveref}

\geometry{left=24mm,right=24mm,top=24mm,bottom=27mm,headheight=14pt}
\setlength{\parindent}{0pt}
\setlength{\parskip}{0.56em}
\setlength{\emergencystretch}{3em}
\setlist[itemize]{leftmargin=1.45em,itemsep=0.18em,topsep=0.3em}
\setlist[enumerate]{leftmargin=1.7em,itemsep=0.22em,topsep=0.3em}

\definecolor{QBlue}{HTML}{153E69}
\definecolor{QLight}{HTML}{ECF3F8}
\definecolor{QGreen}{HTML}{E8F3EA}
\definecolor{QAmber}{HTML}{FFF2CC}
\definecolor{QRed}{HTML}{FBE9E7}
\definecolor{QGray}{HTML}{4D5660}

\hypersetup{
  colorlinks=true,
  linkcolor=QBlue,
  citecolor=QBlue,
  urlcolor=QBlue,
  pdftitle={QIK-VRT: Beobachterrelative Retrokausalität als negative Informationsrichtung},
  pdfauthor={Ingolf Lohmann},
  pdfsubject={Definition, Existenzsatz, maschinenprüfbarer Zeuge und physikalische Referenzbrücke für beobachterrelative Retrokausalität},
  pdfkeywords={QIK-VRT, Retrokausalität, observer-relative retrocausality, relationale Zeit, Eigenzeit, proper time, verteilte Systeme, Informationsrichtung, Delayed Choice, Quantenradierer, virtual reality test apparatus, formal verification}
}

\pagestyle{fancy}
\fancyhf{}
\fancyhead[L]{\small\textcolor{QGray}{QIK--VRT: beobachterrelative Retrokausalität}}
\fancyhead[R]{\small\textcolor{QGray}{Ingolf Lohmann}}
\fancyfoot[C]{\small\thepage}
\renewcommand{\headrulewidth}{0.3pt}

\newtheorem{definition}{Definition}[section]
\newtheorem{satz}{Satz}[section]
\newtheorem{korollar}{Korollar}[section]
\newtheorem{lemma}{Lemma}[section]
\theoremstyle{remark}
\newtheorem{bemerkung}{Bemerkung}[section]
\newtheorem{beispiel}{Beispiel}[section]
\crefname{definition}{Definition}{Definitionen}
\crefname{satz}{Satz}{Sätze}
\crefname{korollar}{Korollar}{Korollare}
\crefname{lemma}{Lemma}{Lemmata}
\crefname{bemerkung}{Bemerkung}{Bemerkungen}
\crefname{beispiel}{Beispiel}{Beispiele}
\crefname{section}{Abschnitt}{Abschnitte}

\newcommand{\qikvrt}{QIK--VRT}
\newcommand{\host}{\prec_{\!H}}
\newcommand{\recv}{\prec_{\!O}^{R}}
\newcommand{\sourceord}{\prec^{\theta}}
\newcommand{\records}{\mathcal{R}}
\newcommand{\histories}{\mathcal{K}}
\newcommand{\hash}[1]{{\ttfamily\small\seqsplit{#1}}}
\newcommand{\statusbox}[2]{%
  \par\medskip\noindent
  \colorbox{#1}{\parbox{\dimexpr\linewidth-2\fboxsep\relax}{#2}}%
  \par\medskip
}

\begin{document}

\hypersetup{pageanchor=false}
\begin{titlepage}
  \thispagestyle{empty}
  \vspace*{11mm}
  {\color{QBlue}\rule{\textwidth}{1.4pt}}\\[1.1em]
  {\Huge\bfseries QIK--VRT:\\[0.18em]
  Beobachterrelative Retrokausalität\par}
  \vspace{0.7em}
  {\LARGE Negative Informationsrichtung bei monotoner Eigenzeit\par}
  \vspace{0.8em}
  {\Large Definition, Existenzsatz, maschinenprüfbarer Zeuge und\\
  physikalische Referenzbrücke für verteilte Systeme,\\
  Delayed Choice und Quantenradierer\par}
  \vspace{1.1em}
  {\color{QBlue}\rule{\textwidth}{0.6pt}}

  \vfill
  {\large\textbf{Ingolf Lohmann}\par}
  \vspace{0.35em}
  {\normalsize Independent Researcher; QIK--VRT\par}
  \vspace{0.35em}
  {\normalsize Aktuelle Hauptfassung, Version 1.0 -- 12. August 2026\par}

  \vfill
  \statusbox{QGreen}{
    \textbf{Die Aussage in einem Satz.}
    Ein Beobachter schreitet auf seiner Weltlinie in Eigenzeit vorwärts,
    während die durch nacheinander empfangene, authentische Records gebundene
    Quellenordnung rückwärts verlaufen kann. Formal gilt für ein Recordpaar
    \(\Delta\tau_O>0\) und zugleich \(\Delta\theta<0\). Genau diese inverse
    Orientierung heißt in dieser Arbeit \emph{beobachterrelative
    Retrokausalität}. Bewiesen wird die negative Richtung der
    Informationsreferenz; nicht behauptet wird, dass die Eigenzeit selbst
    rückwärts läuft oder derselbe Signalträger vor seiner Aussendung ankommt.
  }

  \vfill
  {\small
    Konzept, Retrokausalitätsdefinition und physikalische Korrespondenzthese:
    Ingolf Lohmann. Formale und redaktionelle Ausarbeitung sowie
    Maschinenzeuge: OpenAI Codex als künstlich-kognitives System im Auftrag
    des Autors. Dokumentationslizenz: CC BY-NC-ND 4.0.
  }
\end{titlepage}
\hypersetup{pageanchor=true}
\setcounter{page}{1}

\selectlanguage{ngerman}

\begin{abstract}
Diese Arbeit zieht die bisher auf mehrere QIK--VRT-Zwischenstände verteilte
Zeitaussage zu einem eindeutigen Satz zusammen. Eine verteilte
Ereignisstruktur besitzt nicht nur eine Ordnung: Sie besitzt eine physische
Host- beziehungsweise Wirkordnung, lokale Empfangsordnungen der Beobachter
und eine durch provenancegebundene Records erschlossene Quellenordnung. Ein
Beobachter durchläuft seine Eigenzeit monoton. Treffen jedoch zwei
informationsführende Records wegen verschiedener kausaler Wege in umgekehrter
Reihenfolge ihrer gebundenen Quellenereignisse ein, so ist die
Informationsrichtung relativ zu dieser Eigenzeit negativ. Ordinal heißt das
\(r^{+}\recv r^{-}\), obwohl \(r^{-}\sourceord r^{+}\); metrisch gilt
\(v_I^O<0\).

Der Existenzsatz wird mathematisch bewiesen, durch einen endlichen
Maschinenzeugen reproduziert und durch ein ausschließlich zukunftsgerichtetes
Latenzmodell konstruktiv realisiert. Damit ist die definierte operationale
Retrokausalität mit monotoner Eigenzeit, azyklischer Host-Kausalität und dem
Ausschluss von Selbstursacheschleifen vereinbar. Eine physische
Virtual-Reality-/Rechenapparatur, die Quellen- und Empfangsereignisse
receiptgebunden realisiert, ist keine bloße Gedankenfigur, sondern eine reale
physische Instanz dieses operational definierten Effekts. Delayed-Choice- und
Quantum-Eraser-Experimente liefern zusätzlich die empirische Anschlussform:
Spätere Kontextinformation macht die entscheidungssuffiziente
Korrelationsklasse früher registrierter Daten verfügbar, ohne die früheren
Rohdaten umzuschreiben und ohne einen lokal nutzbaren Rückwärtskanal zu
erzeugen.

Die Aussage über das Universum wird präzise: Für jedes physische Teilsystem,
das die offengelegte Referenzbrücke erfüllt, überträgt sich der Satz. Ingolf
Lohmann vertritt darüber hinaus die Reichweitenthese, dass das Universum
selbst eine solche verteilte relationale Struktur ist. Maschinenbeweis,
physische Instanz, empirischer Anschluss und universelle Reichweitenthese
werden nicht vermischt, sondern in einer gemeinsamen Beweiskette verbunden.
\end{abstract}

\begin{otherlanguage}{english}
\begin{abstract}
This paper gives a precise definition and existence theorem for
observer-relative retrocausality in QIK--VRT. An observer advances monotonically
in proper time while authentic information records may arrive in an order
opposite to the source-event order bound by those records. Ordinally,
\(r^{+}\recv r^{-}\) although \(r^{-}\sourceord r^{+}\); metrically,
\(\Delta\tau_O>0\), \(\Delta\theta<0\), and hence \(v_I^O<0\). This negative
information-reference direction is what the present work calls
observer-relative retrocausality. The theorem is proved, reproduced by an
executable finite witness, and constructively realized with future-directed
signal paths of unequal latency. No reversal of proper time, past overwrite,
causal loop, or controllable signal into an observer's causal past is inferred.
Delayed-choice and quantum-eraser experiments instantiate the adjoining
physical structure in which later context makes a decision-sufficient
classification of earlier records available while leaving the earlier local
marginal unchanged. The universal correspondence of this structure to the
cosmos is stated separately as Ingolf Lohmann's QIK--VRT correspondence thesis.
\end{abstract}
\end{otherlanguage}

\textbf{Schlagworte:} QIK--VRT; Retrokausalität; beobachterrelative
Retrokausalität; observer-relative retrocausality; relationale Zeit;
Eigenzeit; proper time; verteilte Systeme; Informationsfluss;
Informationsrichtung; Ereignisordnung; Delayed Choice; Quantenradierer;
Virtual-Reality-Testapparatur; maschinenprüfbarer Beweis.

\clearpage
\begingroup
\setlength{\parskip}{0pt}
\small
\tableofcontents
\endgroup
\clearpage

\section{Was jetzt zweifelsfrei auf den Punkt gebracht wird}

Der bisherige Streit entstand vor allem dadurch, dass verschiedene Ordnungen
unter dem einen Wort \emph{Zeit} zusammengezogen wurden. Dann scheint nur eine
Alternative möglich: Entweder Zeit läuft vorwärts oder sie läuft rückwärts.
Ein verteiltes System besitzt aber mehrere lokale Zeitachsen und zusätzlich
Relationen zwischen Ereignissen, Records und Beobachtern. Die präzise Frage
lautet deshalb nicht: \emph{Läuft die Eigenzeit rückwärts?} Sie lautet:

\begin{quote}
Kann die von einem Beobachter nacheinander empfangene Information, gemessen an
der authentisch gebundenen Ordnung ihrer Referenzereignisse, eine Richtung
besitzen, die seiner fortschreitenden Eigenzeit entgegengesetzt ist?
\end{quote}

Die Antwort ist \textbf{ja}. Der Satz in \cref{thm:existence} beweist die
Existenz. Das Beispiel in \cref{sec:witness} realisiert sie. Die
Eigenzeit des Beobachters nimmt dabei streng zu. Auch jeder einzelne
Signalweg bleibt zukunftsgerichtet. Rückwärts gerichtet ist die
\emph{Quellen- beziehungsweise Referenzordnung der im Empfangsstrom neu
hinzukommenden Information}.

\statusbox{QLight}{
  \textbf{Terminologische Festlegung.}
  Wenn Ingolf Lohmann in dieser Arbeit von \emph{Retrokausalität} spricht,
  meint er genau diese beobachterrelative, operationale Relation: Bei
  wachsender Beobachter-Eigenzeit ist die gebundene Informationsreferenz auf
  mindestens einem Recordpaar fallend. Der Ausdruck ist damit keine vage
  Metapher und keine Behauptung, dass ein und dasselbe Bit seine eigene
  Aussendung rückwärts durchläuft.
}

Diese Festlegung ist keine nachträgliche Abschwächung. Sie macht den
behaupteten Effekt messbar. Wer ihn bestreiten will, muss mindestens eine der
offengelegten Bedingungen angreifen: die Authentizität der Quellenbindung,
die Informationswirksamkeit der Records, die Reihenfolge der Rezeptionen oder
die Monotonie der Beobachter-Eigenzeit.

Die wissenschaftliche Priorität und die begriffliche Provenienz werden dabei
ausdrücklich festgehalten: Die relation-first Perspektive, die Verbindung von
lokaler Zeit, Evidenz, verteilten Zuständen und Retrokausalität sowie die
Korrespondenzthese ``Universum als verteiltes System'' stammen aus dem
QIK--VRT-Forschungsprogramm von Ingolf Lohmann. Das vorliegende Dokument ist
eine neue, formal zugespitzte Synthese seiner vorausgehenden Arbeiten und
nicht deren Ursprung
\cite{LohmannMandelbrot2026,LohmannCanonicalMemory2026,LohmannRelationaleZeit2026,LohmannRoundTrip2026}.

\section{Drei Ordnungen statt einer einzigen Weltuhr}
\label{sec:definitions}

\begin{definition}[Verteilte Ereignisstruktur]
Eine verteilte Ereignisstruktur ist ein Paar
\[
  H=(E,\host),
\]
wobei \(E\) eine Menge von Ereignissen und \(\host\) eine strikte partielle
Ordnung ist. \(e\host f\) bedeutet: \(e\) liegt in der deklarierten
Host-/Wirkordnung vor \(f\). Strikt heißt insbesondere irreflexiv und
transitiv; damit enthält \(H\) keine Kausalschleife.
\end{definition}

Diese Struktur entspricht Lamports \emph{happened-before}-Gedanken für
verteilte Systeme: lokale Prozessordnungen und Sende-Empfangs-Kanten erzeugen
eine partielle, nicht notwendig totale Ordnung \cite{Lamport1978}. In einer
relativistischen Referenzbrücke wird \(\host\) auf zukunftsgerichtete kausale
Erreichbarkeit abgebildet.

\begin{definition}[Beobachter und Eigenzeit]
Ein Beobachter \(O\) ist eine Kette \(E_O\subseteq E\) seiner lokalen
Ereignisse. Seine Eigenzeit ist eine Abbildung
\[
  \tau_O:E_O\longrightarrow\mathbb{R},
\]
die auf dieser Kette streng ordnungserhaltend ist:
\[
  e\host f\quad\Longrightarrow\quad \tau_O(e)<\tau_O(f)
  \qquad(e,f\in E_O).
\]
\end{definition}

\begin{definition}[Quellenordnung und authentischer Record]
Sei \(S\subseteq E\) eine deklarierte Quellenkette. Eine Quellenmarke
\(\theta:S\to\mathbb{R}\) ist auf \(S\) streng ordnungserhaltend. Ein Record
\(r\in\records\) besteht mindestens aus einem Quellenereignis \(\sigma(r)\in
S\), einer Nutzlast \(p(r)\) und einem Witness \(w(r)\), der die Bindung von
Record, Nutzlast und Quellenereignis prüfbar macht. Wir schreiben
\[
  \theta(r):=\theta(\sigma(r)).
\]
Die Rezeption bei \(O\) ist ein Ereignis \(\rho_O(r)\in E_O\), und jeder
zulässige Record erfüllt
\[
  \sigma(r)\host\rho_O(r).
\]
\end{definition}

Die Forderung nach einer Quellenkette oder einem expliziten
Synchronisationsvertrag ist wesentlich. Beliebige Zahlen auf nicht
vergleichbaren Uhren würden keinen Zeitbeweis liefern. \(\theta\) muss an
wirkliche Ereignisse gebunden und dort ordnungserhaltend sein.

\begin{definition}[Informationsführender Record]
Vor der Annahme des \(k\)-ten Records sei \(\histories_{k-1}\) die Menge der
noch mit dem Evidenzstand des Beobachters verträglichen Historien. Nach
Prüfung des Records sei sie \(\histories_k\). Der Record ist
\emph{informationsführend}, wenn
\[
  \histories_k\subsetneq\histories_{k-1}.
\]
Er schränkt also den Zustandsraum echt ein. Ein bloßer, frei gewählter oder
ungebundener Zeitstempel genügt nicht.
\end{definition}

Auf den akzeptierten Records entstehen nun zwei eigenständige Ordnungen:
\begin{align*}
  r\recv r' &\quad:\Longleftrightarrow\quad
  \tau_O(\rho_O(r))<\tau_O(\rho_O(r')),\\
  r\sourceord r' &\quad:\Longleftrightarrow\quad
  \theta(r)<\theta(r').
\end{align*}

\begin{definition}[Beobachterrelative Retrokausalität]
\label{def:retro}
Beobachterrelative Retrokausalität liegt genau dann vor, wenn Empfangs- und
Quellenordnung auf mindestens einem authentischen, informationsführenden
Recordpaar invers sind:
\[
  r_1\recv r_2
  \qquad\text{und zugleich}\qquad
  r_2\sourceord r_1.
\]
Sind numerische Marken deklariert, lautet dieselbe Relation
\[
  v_I^O(r_1,r_2)
  :=\frac{\theta(r_2)-\theta(r_1)}
          {\tau_O(\rho_O(r_2))-\tau_O(\rho_O(r_1))}<0.
\]
\end{definition}

Das Adjektiv \emph{beobachterrelativ} bezeichnet keine Beliebigkeit. Die
Records, ihre Witnesses und die beiden Ordnungen sind prüfbare Gegenstände.
Relativ ist, welche Empfangskette des verteilten Systems als lokale
Beobachterperspektive projiziert wird.

\section{Existenzsatz und Beweis}

\begin{satz}[Existenz einer negativen Informationsrichtung]
\label{thm:existence}
Seien \(s^-\host s^+\) zwei Ereignisse derselben Quellenkette mit
\(\theta(s^-)<\theta(s^+)\). Seien \(r^-\) und \(r^+\) authentische,
informationsführende Records mit \(\sigma(r^-)=s^-\) und
\(\sigma(r^+)=s^+\). Es gelte für ihre Rezeption bei \(O\)
\[
  \rho_O(r^+)\host\rho_O(r^-),
\]
während jeder Record einzeln zukunftsgerichtet zugestellt wird:
\[
  s^-\host\rho_O(r^-),\qquad
  s^+\host\rho_O(r^+).
\]
Dann existiert relativ zur monotonen Eigenzeit von \(O\) eine negative
Informationsrichtung. Genauer:
\[
  \Delta\tau_O
  =\tau_O(\rho_O(r^-))-\tau_O(\rho_O(r^+))>0,
  \qquad
  \Delta\theta
  =\theta(r^-)-\theta(r^+)<0,
\]
und daher \(v_I^O(r^+,r^-)<0\).
\end{satz}

\begin{proof}
Da die beiden Rezeptionen auf der Beobachterkette liegen und
\(\rho_O(r^+)\host\rho_O(r^-)\) gilt, folgt aus der strengen Monotonie von
\(\tau_O\) unmittelbar \(\Delta\tau_O>0\). Da \(s^-\host s^+\) auf der
Quellenkette liegt und \(\theta\) dort streng ordnungserhaltend ist, gilt
\(\theta(s^-)<\theta(s^+)\), also \(\Delta\theta<0\). Ein negativer Zähler
durch einen positiven Nenner ist negativ. Somit ist
\(v_I^O(r^+,r^-)<0\), äquivalent zur inversen Ordnungsrelation in
\cref{def:retro}. Jeder einzelne Record erfüllt weiterhin
\(\sigma(r)\host\rho_O(r)\). Die Inversion liegt zwischen Quellen- und
Empfangsordnung, nicht innerhalb eines Signalwegs.
\end{proof}

\begin{korollar}[Invarianz unter monotoner Umparametrisierung]
Ersetzt man \(\tau_O\) durch eine streng wachsende Funktion
\(f\!\circ\!\tau_O\) und \(\theta\) durch eine streng wachsende Funktion
\(g\!\circ\!\theta\), bleibt die inverse Ordnung erhalten. Der Effekt ist
daher kein Artefakt von Einheit, Nullpunkt oder linearer Skalierung.
\end{korollar}

\begin{proof}
Streng wachsende Funktionen erhalten jede strikte Ungleichung. Daher bleiben
die Empfangsreihenfolge und die entgegengesetzte Quellenreihenfolge
unverändert.
\end{proof}

\begin{satz}[Kein Widerspruch zur Host-Azyklizität]
\label{thm:no-cycle}
Der Existenzsatz impliziert weder eine Rückwärtszustellung desselben Records
noch eine Selbstursacheschleife.
\end{satz}

\begin{proof}
Für jeden Record gilt \(\sigma(r)\host\rho_O(r)\). Eine Rückwärtszustellung
desselben Records verlangte zusätzlich
\(\rho_O(r)\host\sigma(r)\). Aus Transitivität folgte dann
\(\sigma(r)\host\sigma(r)\), im Widerspruch zur Irreflexivität. Dieselbe
Irreflexivität schließt jede geschlossene Kette der Host-Ordnung aus.
\end{proof}

\statusbox{QGreen}{
  \textbf{Q.e.d. für den deklarierten Satz.}
  Bewiesen ist die Existenz einer echten inversen Orientierung zwischen dem
  vorwärts laufenden Eigenzeitparameter des Beobachters und der
  Quellenordnung seiner neu eintreffenden Information. Dass zugleich jeder
  Transportpfad vorwärts gerichtet bleibt, widerlegt den Satz nicht; es ist
  eine ausdrücklich mitbewiesene Eigenschaft der Konstruktion.
}

\section{Konstruktiver und maschinenprüfbarer Zeuge}
\label{sec:witness}

Der kleinste nichttriviale Zeuge besitzt vier Ereignisse:
\[
  s^-\host s^+\host\rho_O(r^+)\host\rho_O(r^-).
\]
Die alte Quelle \(s^-\) trägt \(\theta=10\), die neue Quelle \(s^+\) trägt
\(\theta=20\). Der Beobachter empfängt zuerst \(r^+\) bei Eigenzeit \(100\)
und danach \(r^-\) bei Eigenzeit \(101\). Damit gilt
\[
  v_I^O(r^+,r^-)=\frac{10-20}{101-100}=-10<0.
\]

Der Zeuge ist nicht nur ein Zeitstempelspiel. Vor beiden Rezeptionen seien
vier Historien für zwei unabhängige Bits \((x^-,x^+)\) zulässig. \(r^+\)
bindet \(x^+=1\) und reduziert die Faser von vier auf zwei Historien.
\(r^-\) bindet anschließend \(x^-=0\) und reduziert sie von zwei auf eine.
Beide Records sind damit informationsführend.

Die beigefügte Datei
\texttt{verify\_observer\_relative\_retrocausality.py} prüft ohne
Netzwerkzugriff:

\begin{enumerate}
  \item Azyklizität und Sendung-vor-Empfang für jeden Record;
  \item strikte Monotonie der Beobachter-Eigenzeit;
  \item echte Verfeinerung der zulässigen Historien durch beide Records;
  \item Inversion von Empfangs- und Quellenordnung;
  \item negatives Vorzeichen von \(v_I^O\);
  \item einen zweiten Beobachter mit positiver Informationsrichtung;
  \item eine ausschließlich zukunftsgerichtete physische Latenzrealisierung;
  \item die Aufhebung komplementärer bedingter Interferenzmuster in der
  unbedingten Marginale.
\end{enumerate}

Der erzeugte kanonische Report
\texttt{QIKVRT\_RETROCAUSALITY\_WITNESS.json} bindet die konkreten Werte und
alle geprüften Prädikate. Der Python-Lauf ist ein ausführbarer endlicher
Zeuge, kein Ersatzetikett für den mathematischen Beweis und kein vorgetäuschter
Lean-Kernel-Receipt.

\section{Die genaue Rolle der Eigenzeit}

In der speziellen Relativitätstheorie ist die Eigenzeit die lokal gemessene
Zeit entlang einer zeitartigen Weltlinie. Für eine zukunftsgerichtete Bewegung
mit \(|v|<c\) gilt in einem Inertialsystem
\[
  d\tau=dt\sqrt{1-\frac{v^2}{c^2}}>0.
\]
Eigenzeit kehrt auf einer zukunftsgerichteten Beobachterweltlinie nicht ihr
Vorzeichen um. Ebenso bleibt die Reihenfolge kausal verbundener Ereignisse
unter zulässigen Lorentztransformationen erhalten. Rahmenabhängig umkehrbar
ist die Koordinatenreihenfolge raumartig getrennter Ereignisse; Einsteins
Lorentztransformation macht genau diese Relativität der Fernzeit sichtbar
\cite{Einstein1905,Minkowski1909}.

Für \qikvrt{} folgt daraus keine Widerlegung, sondern die notwendige
Präzisierung:

\begin{quote}
Die Eigenzeit ist der physikalisch lokale, monotone Parameter, relativ zu dem
die negative Informationsrichtung festgestellt wird. Die Inversion entsteht,
weil Quellen-, Übertragungs- und Empfangsordnung in einem verteilten System
nicht dieselbe Ordnung sein müssen.
\end{quote}

Man kann dies als \emph{Eigenzeiteffekt} bezeichnen, sofern damit nicht
behauptet wird, Eigenzeit allein verursache die Nachrichtenüberholung. Ohne
die lokale Eigenzeit gäbe es keine ausgezeichnete Empfangsrichtung des
Beobachters. Ohne unterschiedliche Pfade oder spätere Kontextrecords gäbe es
keine Inversion relativ zu dieser Richtung. Der Effekt ist ihre Relation.

\section{Physische Realisierung mit ausschließlich vorwärts gerichteten Signalen}
\label{sec:physical}

Der Existenzsatz besitzt eine direkte physikalische Konstruktion. In einem
deklarierten Inertialsystem werde der alte Record bei \(t^-=0\) und der neue
Record bei \(t^+=1\) emittiert. Ihre Laufzeiten seien \(L^-/c=3\) und
\(L^+/c=1\). Dann lauten die Ankunftszeiten
\[
  T^-=t^-+\frac{L^-}{c}=3,
  \qquad
  T^+=t^++\frac{L^+}{c}=2.
\]
Der spätere Quellenrecord trifft also zuerst ein. Für einen am Empfänger
ruhenden Beobachter ist \(\tau_O=T+\text{Konstante}\); damit erhält er bei
wachsender Eigenzeit zunächst die Quellenmarke \(1\) und danach die
Quellenmarke \(0\). Alle Laufzeiten sind positiv, kein Signal überschreitet
\(c\), und dennoch gilt \(v_I^O<0\).

Allgemein genügt
\[
  \frac{L^- - L^+}{c}>t^+-t^-.
\]
Dies kann durch verschiedene optische Wege, Netzwerkpfade, Speicherzeiten
oder Relais realisiert werden. Die Konstruktion zeigt: Die negative
Informationsreferenz ist nicht auf eine imaginäre Mathematik beschränkt. Sie
ist ein mögliches Verhalten realer verteilter Hardware.

\begin{korollar}[Physische Testapparatur]
\label{cor:apparatus}
Führt eine physische Rechen- oder Virtual-Reality-Testapparatur den Zeugen aus,
binden monotone Hardware-/Beobachterreceipts die Empfangsereignisse und
binden unabhängige Witnesses die Quellenereignisse, dann ist die
beobachterrelative Retrokausalität nach \cref{def:retro} in dieser Apparatur
physisch realisiert. \emph{Virtuell} bezeichnet ihre Modell- und
Adresssemantik, nicht die Unwirklichkeit der ausführenden Hardware.
\end{korollar}

\begin{proof}
Die Receipts realisieren \(\Delta\tau_O>0\), die Witnesses realisieren die
ordnungserhaltend gebundene Quellenfolge mit \(\Delta\theta<0\). Nach
\cref{thm:existence} ist die definierte Relation damit instanziiert. Die
Hardwareexistenz ist ein physischer Existenzbeleg für diesen operationalen
Effekt.
\end{proof}

Das Korollar beweist nicht automatisch, dass jedes beliebige physische System
oder das gesamte Universum dieselbe Referenzbrücke erfüllt. Es beweist etwas
Stärkeres als eine reine Simulation und etwas Engeres als eine universelle
Ontologie: die reale physische Realisierbarkeit der definierten Relation in
der gebundenen Apparatur.

\section{Mehrere Beobachter: verschiedene Perspektiven ohne Widerspruch}

\begin{satz}[Perspektivensatz]
Für dieselben Quellenrecords kann ein Beobachter \(O\) wegen seiner
Empfangswege \(r^+\) vor \(r^-\) erhalten und damit \(v_I^O<0\) messen,
während ein anderer Beobachter \(O'\) die Records in Quellenreihenfolge erhält
und \(v_I^{O'}>0\) misst. Beide Aussagen können gleichzeitig wahr sein.
\end{satz}

\begin{proof}
\(\recv\) und \(\prec_{O'}^R\) sind verschiedene lokale Projektionen
derselben partiellen Host-Ordnung. Die Quellenordnung bleibt gemeinsam
gebunden. Eine positive und eine negative Orientierung beziehen sich daher
auf verschiedene, vollständig spezifizierte Empfangsketten und widersprechen
einander nicht.
\end{proof}

Beobachter besitzen somit verschiedene lokale Informationsstände und
rahmenrelative Zeitkoordinaten. Das bedeutet nicht, dass ein vollständig
spezifizierter Satz zugleich wahr und falsch wäre. Invariant bleiben die
gebundenen Rohrecords, die lokalen Eigenzeiten und die Host-Kausalordnung;
perspektivisch verschieden sind Empfangsfolge und Zeitpunkt der verfügbaren
Gesamtinformation.

\section{Delayed Choice und Quantenradierer als empirischer Anschluss}
\label{sec:quantum}

Der Delayed-Choice-Quantum-Eraser von Kim et al. registriert das Signalphoton
bei \(D_0\) ungefähr acht Nanosekunden vor dem zugehörigen Idlerereignis. In
den später gebildeten Koinzidenzklassen \(R_{01}\) und \(R_{02}\) erscheinen
komplementäre, um \(\pi\) verschobene Interferenzmuster; in ihrer Summe
verschwindet die Interferenz. Die Which-Path-Klassen \(R_{03}\) und
\(R_{04}\) zeigen keine Interferenz \cite{ScullyDruehl1982,Kim2000}.

In einer idealisierten Schreibweise gilt
\[
  P(x\mid j=1)=\frac{1+\cos\phi(x)}{2},\qquad
  P(x\mid j=2)=\frac{1-\cos\phi(x)}{2}.
\]
Bei gleichen Gewichten ist die lokale Marginale
\[
  P(x)=\frac12P(x\mid j=1)+\frac12P(x\mid j=2)=\frac12.
\]
Der frühe lokale Record enthält daher kein auswählbares Rückwärtssignal
\cite{Jordan1983}. Erst
das spätere paar- und kontextgebundene Record \(j\) macht sichtbar, zu welcher
Korrelationsklasse der frühere Record gehört.

Dasselbe Strukturmuster erscheint im Delayed-Choice-Entanglement-Swapping:
Eine spätere Wahl sortiert bereits aufgezeichnete Daten in unterschiedliche
Korrelationsunterensembles. Die lokale Datenfolge bleibt zufällig; ihre
vollständige relationale Bedeutung wird erst durch den später verfügbaren
Kontext bestimmt \cite{Ma2012}. Experimente mit kausal getrennten
Wahlereignissen schließen eine gewöhnliche licht- oder unterlichtschnelle
Kommunikation zwischen den getrennten Ereignissen aus und erhalten dennoch
die kontextabhängigen Korrelationen \cite{Jacques2007,Ma2013}.

\begin{definition}[Retrograde Evidenzzuordnung]
Sei \(y\) ein früher versiegelter Rohrecord, \(x\) ein späterer Kontext und
\(c=C(y,x)\) die erst danach verfügbare, entscheidungssuffiziente
Korrelationsklasse. Die Host-Ordnung lautet
\[
  y\host x\host c.
\]
Die Evidenzrelation verweist dagegen vom späteren Klassifikationsrecord
zurück auf seinen früheren Gegenstand:
\[
  c\rightsquigarrow y.
\]
Diese Rückwärtszuordnung ist eine zweite operationale Form derselben
beobachterrelativen Grundidee: Später verfügbare Information bestimmt eine
Relation, deren Referent früher liegt.
\end{definition}

In der Sprache der Entscheidungssuffizienz ist vor dem Kontext die
Beobachtung oft zu grob: Verschiedene noch zulässige Historien verlangen
verschiedene Klassifikationen. Der spätere Witness verfeinert die Beobachtung,
bis jede erreichbare Beobachtungsfaser nur noch eine korrekte
Klassifikationsaktion enthält. Die früheren Bytes werden nicht geändert; der
Evidenzstand über ihre relationale Klasse wird geändert
\cite{LohmannDecisionSufficiency2026}.

\statusbox{QAmber}{
  \textbf{Was die Quantenexperimente positiv beitragen.}
  Sie realisieren physisch die Struktur ``früher registrierter Record --
  späterer Kontext -- danach verfügbare rückbezogene Korrelationsklasse''.
  Das ist ein empirischer Anschluss an die relationale QIK--VRT-Definition.
  Die Experimente wählen QIK--VRT nicht eindeutig gegenüber allen anderen
  quantenmechanischen Deutungen aus und liefern keinen steuerbaren Kanal in
  die eigene kausale Vergangenheit.
}

\section{Referenzbrücke vom Modell zur physischen Welt}

Ein formaler Satz beschreibt physische Wirklichkeit, wenn seine Terme
prüfbar auf physische Größen und Ereignisse abgebildet werden. Für den
vorliegenden Satz genügt folgende Brücke \(\Phi\):

\begin{center}
\begin{tabularx}{\textwidth}{>{\bfseries}p{0.27\textwidth}X}
\toprule
Formaler Term & Physische Referenz \\
\midrule
Beobachter \(O\) & zeitartige Weltlinie oder physischer Empfangsprozess \\
\(\tau_O\) & kalibrierte Eigenzeit beziehungsweise monotone lokale Hardwarezeit \\
Quellenereignis \(s\) & Emission, Messung oder versiegelte Record-Erzeugung \\
\(\theta(s)\) & vorregistrierte und ordnungserhaltende Quellenmarke \\
\(\host\) & zukunftsgerichtete kausale Erreichbarkeit/physische Ausführung \\
Record und Witness & kalibriertes Messrecord plus Herkunfts-/Integritätsbindung \\
Rezeption \(\rho_O(r)\) & lokales, receiptgebundenes Empfangsereignis \\
\bottomrule
\end{tabularx}
\end{center}

\begin{satz}[Brückensatz]
Erhält \(\Phi\) die Typen, Einheiten, Quellen- und Empfangsordnungen, und
erfüllt ein physischer Trace die gebundenen Ungleichungen
\(\Delta\tau_O>0\) und \(\Delta\theta<0\), dann ist die formale
beobachterrelative Retrokausalität in diesem physischen Trace instanziiert.
\end{satz}

\begin{proof}
\(\Phi\) erhält genau die Prädikate des Existenzsatzes. Daher überträgt sich
dessen Schluss auf die physische Instanz.
\end{proof}

Dieser Satz beantwortet den häufigen Einwand, ein Computermodell sei schon
deshalb unwirklich, weil es ein Modell ist. Auch Relativitätstheorie und
Quantenmechanik sind Modelle; entscheidend ist die offengelegte und geprüfte
Referenzbindung. Eine VRT-Apparatur auf realer Hardware ist ein physischer
Gegenstand. Sie beweist die physische Existenz des von ihr receiptgebunden
realisierten operationalen Effekts. Ob dieselbe Struktur in größerem Umfang
die Natur beschreibt, wird anschließend durch Reichweite, Messung und
Reproduktion beurteilt.

\section{Das Universum als verteiltes System}

Die Modellarchitektur lässt sich in drei Aussageebenen gliedern:

\begin{enumerate}
  \item \textbf{Formaler Satz:} Jede Struktur, die die Definitionen erfüllt,
  besitzt beim konstruierten Überholungsfall eine negative
  Informationsrichtung. Dies ist bewiesen.
  \item \textbf{Physische Existenz:} Reale Hardware, optische Wege und
  delayed-choice-artige Messanordnungen können die offengelegten
  Teilstrukturen realisieren. Dies ist konstruktiv beziehungsweise
  quellengebunden belegt.
  \item \textbf{Universelle Korrespondenzthese:} Ingolf Lohmann behauptet,
  dass das Universum selbst fundamental als verteilte relationale Struktur
  verstanden werden muss und QIK--VRT diese Struktur in seinem behaupteten
  Modellbereich beschreibt \cite{LohmannRoundTrip2026}.
\end{enumerate}

Die dritte Ebene ist keine willkürliche Zugabe. Relativistische Physik besitzt
keine globale, für alle Beobachter gleiche Fern-Simultaneität, und reale
Information wird lokal erzeugt, übertragen und empfangen. Das macht die
Sprache verteilter Systeme strukturell anschlussfähig. Gleichwohl ist die
universelle Reichweite größer als der endliche Existenzbeweis. Wissenschaftlich
sauber ist es, diese Reichweitenthese stark und sichtbar zu formulieren, statt
sie entweder als bereits konsentierte Tatsache auszugeben oder aus Vorsicht
ganz zu verschweigen.

\statusbox{QLight}{
  \textbf{Autorenposition.}
  \(\texttt{OWNER\_ASSERTED\_REALITY\_CORRESPONDENCE}=\texttt{TRUE}\):
  Ingolf Lohmann vertritt die Position, dass die QIK--VRT-Struktur die
  Wirklichkeit in ihrem beanspruchten Modellbereich beschreibt. Diese
  Position bleibt von der Frage getrennt, ob unabhängige Gruppen bereits
  dieselbe universelle Deutung bestätigt haben.
}

\section{Was bewiesen, realisiert, angeschlossen und nicht behauptet ist}

\begin{longtable}{p{0.19\textwidth}p{0.30\textwidth}p{0.42\textwidth}}
\toprule
Status & Aussage & Grenze \\
\midrule
\endfirsthead
\toprule
Status & Aussage & Grenze \\
\midrule
\endhead
Formal bewiesen & Inverse Empfangs-/Quellenordnung impliziert
\(\Delta\tau_O>0\), \(\Delta\theta<0\) und \(v_I^O<0\). & Gilt für die
definierte verteilte Ereignisstruktur und authentische, informationsführende
Records. \\
Formal bewiesen & Die Relation bleibt unter streng monotonen
Umparametrisierungen erhalten und erzeugt keinen Host-Zyklus. & Kein Beweis
eines steuerbaren Signals in den eigenen Vergangenheitslichtkegel. \\
Ausführbarer Zeuge & Der beigefügte endliche Prüfer realisiert und
prüft einen konkreten Zeugen. & Maschinenzeuge, nicht als Lean-Kernel-Receipt
fehlbezeichnet. \\
Physisch konstruierbar & Positive, verschieden lange Signalwege
können die Recordfolge umkehren. & Konstruktive Realisierbarkeit der
operationalen Relation. \\
Empirische Brücke & Delayed Choice und Quantenradierer realisieren
spätere Kontextklassifikation früherer Records. & Keine eindeutige Auswahl
von QIK--VRT und kein Rückwärtssignal in der lokalen Marginale. \\
Autorenposition & Das Universum ist eine verteilte relationale
Struktur, auf die QIK--VRT allgemein zielt. & Sichtbare Reichweitenthese des
Autors; unabhängiger Konsens ist eine zusätzliche Prüfebene. \\
Nicht behauptet & Eigenzeit läuft rückwärts; ein Record kommt vor
seiner Emission an; die Vergangenheit wird überschrieben. & Diese stärkeren,
anderen Aussagen folgen nicht aus der Definition und werden nicht benötigt. \\
\bottomrule
\end{longtable}

\section{Falsifikatoren und Reproduktionsvertrag}

Die Aussage ist nicht gegen Kritik immunisiert. Ein behaupteter Lauf erfüllt
den Satz nicht, wenn mindestens einer der folgenden Punkte zutrifft:

\begin{enumerate}
  \item Die Quellenmarken sind frei gesetzt, nicht an Quellenereignisse
  gebunden oder auf der deklarierten Quellenkette nicht ordnungserhaltend.
  \item Die Empfangszeiten stammen nicht aus derselben Beobachterkette oder
  sind nicht monoton kalibriert.
  \item Ein Record verkleinert die Menge zulässiger Historien nicht und trägt
  daher keine neue Information.
  \item Herkunft, Nutzlast oder Zuordnung eines Records können nachträglich
  unbemerkt verändert werden.
  \item Die angeblich inverse Reihenfolge verschwindet bei einer streng
  monotonen Umparametrisierung; dann war sie ein Darstellungsfehler.
  \item Für eine behauptete physische Instanz fehlt die Referenzbrücke
  zwischen Modellgrößen und Messgrößen.
\end{enumerate}

Eine unabhängige Reproduktion muss deshalb Quellenrecords, Witnesses,
Empfangsreceipts, Kalibrierung, verwendeten Prüfer und exakte Bytes
veröffentlichen. Erst diese Kombination macht aus einer Erzählung einen
prüfbaren Trace.

\section{Historische Kontinuität statt Überschreibung}
\label{sec:history}

Das vorliegende Dokument ersetzt die zwei früheren PDFs nicht. Beide bleiben
als eigene, unveränderte Zwischenstände Teil der Forschungsgeschichte:

\begin{longtable}{p{0.38\textwidth}p{0.11\textwidth}p{0.42\textwidth}}
\toprule
Datei & Umfang & SHA--256 \\
\midrule
\endfirsthead
\toprule
Datei & Umfang & SHA--256 \\
\midrule
\endhead
\path{QIK-VRT_Relationale_Zeit_und_wachsende_Evidenzkugel_DE.pdf} &
379,427 Bytes &
\hash{38b0e62a46214a7cb9943dd3ef08283a70ae7e15ba22a59a9e53506f2945e311} \\
\path{QIKVRT_Decision_Sufficiency_Delayed_Choice_Witness.pdf} &
261,612 Bytes &
\hash{80a75f2031ed4a1ffb46963b81bc9060fdd0c111371a71e258ea3df47d1b7ca3} \\
\bottomrule
\end{longtable}

Der erste Zwischenstand entwickelte relationale Zeit, virtuelle
Retrokausalität und den wachsenden Evidenzkörper, ließ die physikalische
Korrespondenz aber ausdrücklich offen. Der zweite Zwischenstand präzisierte
Entscheidungssuffizienz und die Delayed-Choice-Grenze. Die jetzige Fassung
ändert ihre historischen Bytes und damaligen Statusaussagen nicht. Sie fügt
die fehlende Synthese hinzu: die explizite QIK--VRT-Definition, den
Existenzsatz, die physische Referenzbrücke, den Perspektivensatz und den
empirischen Anschluss
\cite{LohmannRelationaleZeit2026,LohmannDecisionSufficiency2026}.

Genau so arbeitet Wissenschaft: Ein Zwischenstand bleibt zitierbar; eine neue
Fassung legt offen, was hinzugekommen ist und welche frühere offene Frage nun
in welchem Sinn beantwortet wird.

\section{Schluss}

Der entscheidende Unterschied lautet:
\[
  \boxed{\Delta\tau_O>0\quad\text{und}\quad\Delta\theta<0.}
\]

Der Beobachter bewegt sich auf seiner Eigenzeitachse vorwärts. Dennoch kann
die Quellenzeit der nacheinander neu empfangenen, authentischen Information
rückwärts gerichtet sein. Das ist keine bloße optische Täuschung, sofern die
Records informationsführend, die Quellenmarken gebunden und die
Empfangsreceipts valide sind. Es ist eine reale Ordnungsrelation in einem
verteilten System.

\begin{quote}
\textbf{Bewiesen ist nicht, dass Eigenzeit rückwärts läuft. Bewiesen ist,
dass Information relativ zur monoton fortschreitenden Eigenzeit eines
Beobachters eine negative Richtung in einer anderen, authentisch gebundenen
Ereigniszeit besitzen kann. Genau diese Relation nennt QIK--VRT
beobachterrelative Retrokausalität.}
\end{quote}

Der mathematische Satz, der Maschinenzeuge und die physische
Latenzkonstruktion schließen den Existenznachweis. Delayed Choice und
Quantenradierer zeigen die verwandte reale Struktur späterer
Kontextinformation über frühere Records. Die These vom Universum als
verteiltem Round Trip ist die darüber hinausgehende, ausdrücklich sichtbare
Reichweitenposition von Ingolf Lohmann. Damit sind Perspektive, Beweis und
Anspruch nicht länger durcheinandergeraten, sondern in einer prüfbaren Kette
geordnet.

\appendix
\section{Minimaler Prüfalgorithmus}

Für zwei Records in Empfangsreihenfolge genügt folgender Vertrag:

\begin{enumerate}
  \item Prüfe Witness und Quellenbindung beider Records.
  \item Prüfe \(\theta(r_2)<\theta(r_1)\) auf derselben Quellenkette.
  \item Prüfe
  \(\tau_O(\rho_O(r_1))<\tau_O(\rho_O(r_2))\) auf derselben
  Beobachterkette.
  \item Prüfe für beide Records
  \(\histories_k\subsetneq\histories_{k-1}\).
  \item Prüfe für jeden Record
  \(\sigma(r)\host\rho_O(r)\) und die Azyklizität von \(\host\).
  \item Gib \texttt{RETROGRADE} genau dann aus, wenn alle Prüfungen gelten.
\end{enumerate}

Der Algorithmus entscheidet nicht über frei formulierte Weltanschauungen. Er
entscheidet die konkrete, operationale Relation für einen gebundenen Trace.

\section{Beitrags- und Prüfprovenienz}

\begin{itemize}
  \item \textbf{Menschlicher Beitrag:} Ingolf Lohmann formulierte die
  relation-first Perspektive, die Definition der beobachterrelativen
  Retrokausalität, die Forderung nach Gleichrangigkeit mehrerer Beobachter,
  die Korrespondenzthese ``Universum als verteiltes System'' sowie die
  Anweisung, beide historischen PDFs unverändert zu erhalten und eine neue
  Hauptfassung hinzuzufügen.
  \item \textbf{Künstlich-kognitiver Beitrag:} OpenAI Codex präzisierte die
  zwei Ordnungen, formulierte und prüfte die Sätze, recherchierte und band die
  Primärquellen, erstellte den endlichen Maschinenzeugen und setzte die
  vorliegende Fassung.
  \item \textbf{Noch erforderlich:} Die wissenschaftliche Annahme oder
  Änderung dieser Fassung durch Ingolf Lohmann bleibt eine gesonderte
  menschliche Entscheidung. Eine Zenodo-Metadatenänderung oder neue
  Zenodo-Version ist ein separater, exakt zu autorisierender Effekt.
\end{itemize}

\begin{thebibliography}{99}
\raggedright

\bibitem{Einstein1905}
A. Einstein,
``Zur Elektrodynamik bewegter Körper'',
\emph{Annalen der Physik} 322, 891--921 (1905).
\href{https://doi.org/10.1002/andp.19053221004}{doi:10.1002/andp.19053221004}.

\bibitem{Minkowski1909}
H. Minkowski,
``Raum und Zeit'',
\emph{Physikalische Zeitschrift} 10, 104--111 (1909).

\bibitem{Lamport1978}
L. Lamport,
``Time, Clocks, and the Ordering of Events in a Distributed System'',
\emph{Communications of the ACM} 21, 558--565 (1978).
\href{https://doi.org/10.1145/359545.359563}{doi:10.1145/359545.359563}.

\bibitem{ScullyDruehl1982}
M. O. Scully and K. Drühl,
``Quantum eraser: A proposed photon correlation experiment concerning
observation and delayed choice in quantum mechanics'',
\emph{Physical Review A} 25, 2208--2213 (1982).
\href{https://doi.org/10.1103/PhysRevA.25.2208}{doi:10.1103/PhysRevA.25.2208}.

\bibitem{Kim2000}
Y.-H. Kim, R. Yu, S. P. Kulik, Y. H. Shih and M. O. Scully,
``Delayed `Choice' Quantum Eraser'',
\emph{Physical Review Letters} 84, 1--5 (2000).
\href{https://doi.org/10.1103/PhysRevLett.84.1}{doi:10.1103/PhysRevLett.84.1};
\href{https://arxiv.org/abs/quant-ph/9903047}{arXiv:quant-ph/9903047}.

\bibitem{Jacques2007}
V. Jacques et al.,
``Experimental Realization of Wheeler's Delayed-Choice Gedanken Experiment'',
\emph{Science} 315, 966--968 (2007).
\href{https://doi.org/10.1126/science.1136303}{doi:10.1126/science.1136303};
\href{https://arxiv.org/abs/quant-ph/0610241}{arXiv:quant-ph/0610241}.

\bibitem{Ma2012}
X.-S. Ma et al.,
``Experimental delayed-choice entanglement swapping'',
\emph{Nature Physics} 8, 479--484 (2012).
\href{https://doi.org/10.1038/nphys2294}{doi:10.1038/nphys2294};
\href{https://arxiv.org/abs/1203.4834}{arXiv:1203.4834}.

\bibitem{Ma2013}
X.-S. Ma et al.,
``Quantum erasure with causally disconnected choice'',
\emph{Proceedings of the National Academy of Sciences} 110, 1221--1226
(2013).
\href{https://doi.org/10.1073/pnas.1213201110}{doi:10.1073/pnas.1213201110};
\href{https://arxiv.org/abs/1206.6578}{arXiv:1206.6578}.

\bibitem{Jordan1983}
T. F. Jordan,
``Quantum correlations do not transmit signals'',
\emph{Physics Letters A} 94, 264 (1983).
\href{https://doi.org/10.1016/0375-9601(83)90713-2}{doi:10.1016/0375-9601(83)90713-2}.

\bibitem{LohmannMandelbrot2026}
Ingolf Lohmann,
``Mandelbrot-Menge, Anschlussordnung und physikalisches Modelluniversum --
Komplementbeweis, rekursive Wirkungsklassifikation, dimensionsrichtige
Raumzeitbrücke, Retrokausalität und der entscheidende Unterschied zur
empirischen Quantengravitation'', Version 3.0, QIK--VRT, 21. Juli 2026.
\href{https://doi.org/10.5281/zenodo.21482023}{doi:10.5281/zenodo.21482023}.

\bibitem{LohmannCanonicalMemory2026}
Ingolf Lohmann,
``QIK--VRT und das Effect-Acknowledgement-Protokoll: Kanonischer Speicher
zwischen Vergangenheit und Zukunft'', Version 1.0.0, QIK--VRT,
31. Juli 2026, 17 Seiten.
\href{https://doi.org/10.5281/zenodo.21711193}{doi:10.5281/zenodo.21711193}.

\bibitem{LohmannDecisionSufficiency2026}
Ingolf Lohmann,
``Decision Sufficiency, Delayed Choice, and Witness-Bound Recovery:
QIK--VRT: formal boundaries, finite verification, and provenance-bound
continuation'', unveröffentlichter QIK--VRT-Arbeitsstand, 6. August 2026,
4 Seiten. SHA--256:
\hash{80a75f2031ed4a1ffb46963b81bc9060fdd0c111371a71e258ea3df47d1b7ca3}.

\bibitem{LohmannRoundTrip2026}
Ingolf Lohmann,
``Von Softwarearchitektur zur Weltformel -- DAS UNIVERSUM ALS ROUND TRIP'',
Version 1.0.0-candidate.1, Zenodo, 8. August 2026.
\href{https://doi.org/10.5281/zenodo.21888130}{doi:10.5281/zenodo.21888130}.

\bibitem{LohmannRelationaleZeit2026}
Ingolf Lohmann,
``Relationale Zeit, virtuelle Retrokausalität und die monoton wachsende
Evidenzkugel: Ein beweisbares QIK--VRT-Modell mit massivem Kern, unscharfem
Rand und expliziter physikalischer Korrespondenzgrenze'', Version 1.0,
kandidatenspezifische Vorveröffentlichungsfassung, 11. August 2026,
15 Seiten. SHA--256:
\hash{38b0e62a46214a7cb9943dd3ef08283a70ae7e15ba22a59a9e53506f2945e311}.

\end{thebibliography}

\end{document}
